\section{What was inspected, and where the new work begins}

\subsection{Pinned sources and the proved scope}

The upstream source is \texttt{lsernau/grh-lean} \cite{GRHLean},
pinned at commit
\begin{center}
\texttt{72fead102f663d5c7a407964a13dc0cf9d977207}.
\end{center}
The companion source is the merged article in
\texttt{VladimirReshetnikov/ProveIt} \cite{ProveItArticle}, pinned at commit
\begin{center}
\texttt{22e5c8b7dda558fd9a50b61465483cb4c0433694}.
\end{center}
These references identify the material used here, even if either
repository subsequently changes.

The upstream \texttt{Statements.lean} imports Mathlib alone. Its types
are about Mathlib's Riemann zeta function and Dirichlet $L$-functions.
The five results exported by \texttt{Main.lean} include
\texttt{riemannHypothesis} and
\texttt{LFunction\_eq\_zero\_re}, the latter asserting that every
Dirichlet zero in the open critical strip has real part $1/2$.
The main statements do not assume a substitute analytic inequality as an
extra premise. The scope is RH and GRH for Dirichlet characters; statements
about general Artin or automorphic $L$-functions do not follow merely by
changing the name of the family.

Andrew Sutherland's primary verification record, updated on 27 November
2026, reports successful builds, statement comparisons, axiom audits and
independent kernel checks for both editions at this commit
\cite{Sutherland2026}. The recorded axiom dependencies are
\texttt{propext}, \texttt{Classical.choice} and \texttt{Quot.sound}.
Our work inspected the pinned source, the exported statements and the
main reduction, and performed a source scan. We did not independently
rebuild the upstream Lean development. Published kernel checks and this
source review are different pieces of evidence.

\subsection{The transferable mechanism}

The source's central comparison combines sufficiently many short forms
with a lower bound on the size of integral forms. The determinant and
product formula enforce a schematic estimate
\[
 m\,\mathrm{remainder}\le f a(\log b)^5,
 \qquad m\asymp\sqrt{x}.
\]
Its arithmetic reduction gives a one-sided estimate of the shape
\[
 R_{q,d}(x)\le x^{-1/2+\eta}\quad\text{eventually for every }\eta>0,
\]
for a smoothed prime-power remainder carrying a quadratic-character
weight. The final analytic argument is explicit in
\texttt{New/Reduction/Zeros}: the transformed remainder has residues
given by differences of zero multiplicities, an auxiliary twist avoids
the specified zero region, and a Landau sign argument excludes zeros to
the right of the critical line. Functional equations finish the reduction.
The reading guide and the proof files describe this chain
\cite{GRHLean}.

We use the proved zero-location theorem and its positive spectral
consequences. We do not presume that the source's deep determinant
construction transfers unchanged to a new family. The additional
arguments in this note are harmonic analysis, convexity, algebraic
separation and explicit-formula analysis, proved below.

\subsection{The precise additions to the collection}

The merged article describes itself as a synthesis of 103 reports,
containing 760 results in 621 pages; it also explicitly says that the
merge was not an independent mathematical review \cite{ProveItArticle}.
All its section source files were retrieved, and the relevant inverse
and filter sections were read in detail. The following comparison fixes
the contribution of the present work.

\begin{center}
\small
\begin{tabular}{@{}p{.30\linewidth}p{.31\linewidth}p{.31\linewidth}@{}}
\toprule
Question & Existing result & Addition proved here\\
\midrule
One real nonnegative weight & One complex weight encodes all labelled
Dirichlet zero spectra & One real positive weight does so, with a known
prime-only correction, and may be arbitrarily close to one\\[4pt]
Nonconsecutive zero frequencies & Sharp single-frequency jets and
consecutive-harmonic averages & Every pair of frequencies; all
minimizers; cube-root versus linear resonance asymptotics; exact
shortening slope\\
\bottomrule
\end{tabular}
\end{center}

The arithmetic result uses exact limiting data. It is not a claim that a
finite noisy measurement can recover infinitely many zeros. Its novelty
lies in the positive real prime-side construction and the recovery of
the labelled oriented spectra. The underlying irrational-reflection
mechanism is related to established phase-retrieval results
\cite{Liehr2024,Wellershoff2024}.

The filter problem is also the undamped positive two-mode input-shaping
problem from vibration control. The low-ratio formula, and numerical
descriptions of simultaneous two-mode shapers, have older precedents
\cite{Singhose1997,Crain1996}. We therefore do not claim that positive
two-frequency cancellation was introduced here. The quantitative
resonance theorem and its shortening and stability consequences are
the proposed new refinement. The supplied collection does not contain
these two extensions. A targeted search cannot certify worldwide
historical priority; the proofs stand independently of that question.

\subsection{A reading route}

For the main arithmetic theorem, start with
Theorem~\ref{rw:thm:encoder}, then follow the coefficient formula
\eqref{rw:eq:weight}, the amplitude recovery
\eqref{rw:eq:decodeA}, and the terminating reflection reconstruction
\eqref{rw:eq:Pelltelescope}. For the extremal theorem, start with
Theorems~\ref{filt:classification} and~\ref{filt:cusp}.
The nonnegative function in~\eqref{filt:g} is the certificate that
no shorter positive average can work. Appendix~\ref{app:verification}
describes the reproducible checks and the exact scope of verification.
